\section{Stationary dq machine model}

\begin{assumption}[Steady-state linear model]
The rotor-oriented $dq$ variables are constant, $L_d,L_q>0$, $R_s\geq0$, and
the permanent-magnet flux linkage $\psi_f$ is constant.  Electrical speed is
$\omega_e$.  Motor torque and positive $i_q$ use the same sign convention.
\end{assumption}

The flux linkages and stator voltages are
\begin{align}
 \psid&=\Ld\id+\psif, & \psiq&=\Lq\iq,\label{eq:psm-flux}\\
 \ud&=\Rs\id-\omegae\Lq\iq, &
 \uq&=\Rs\iq+\omegae(\Ld\id+\psif).\label{eq:psm-voltage-model}
\end{align}
The terms proportional to $R_s$ are resistive drops; the terms proportional
to $\omega_e$ are rotational electromotive forces.  The inverter imposes
\begin{equation}
 \ud^2+\uq^2\leq\Umax^2.\label{eq:voltage-ball}
\end{equation}
This is a disk in voltage coordinates.  The task is to determine its inverse
image in current coordinates.

Set
\begin{equation}
 \current=\begin{bmatrix}\id\\\iq\end{bmatrix},\quad
 \voltage=\begin{bmatrix}\ud\\\uq\end{bmatrix},\quad
 \Apsm=\begin{bmatrix}\Rs&-\omegae\Lq\\
             \omegae\Ld&\Rs\end{bmatrix},\quad
 \vect{r}=\begin{bmatrix}0\\\omegae\psif\end{bmatrix}.
\end{equation}
Then the two scalar equations become the affine map
\begin{equation}
 \voltage=\Apsm\current+\vect{r}.\label{eq:affine-map}
\end{equation}

\geointerpretation The linear part rotates and directionally scales current
displacements; $\vect r$ translates their image in voltage space.

\physinterpretation The PM flux produces voltage even at zero stator current.
Resistance and inductive coupling determine how a current increment changes
the two voltage components.

\engineeringconsequence A voltage circle need not remain a circle when pulled
back through the machine map.  Its precise image is governed by the singular
values and invertibility of $\Apsm$.
